% !TEX root = ../notes.tex

\documentclass[../notes.tex]{subfiles}

\begin{document}

\section{March 11}
Here we go.

\subsection{Thick Subcategories of Finite Spectra}
To convince ourselves that the type is a good invariant on our category, we should check that it has good categorical properties.
\begin{theorem} \label{thm:thick-finite-subcat}
	For each $n\ge0$, the category $\mc C_{\ge n}$ of finite spectra of type at least $n$ is a thick subcategory of $\mathbb S_{(p)}$-modules. Conversely, any thick subcategory of finite spectra is $\{0\}$ or takes the form $\mc C_{\ge n}$ for some $n\ge0$.
\end{theorem}
\begin{example}
	It would follow from \Cref{ex:s-eta-type-zero} and \Cref{thm:thick-finite-subcat} that the thick subcategory generated by $\mathbb S/\eta$ should have all finite spectra. Well, $\mathbb S/\eta^2$ is an extension of $\mathbb S/\eta$ by itself, and one can similarly construct $\mathbb S/\eta^3$ and $\mathbb S/\eta^4$. But $\eta^4=0$, so $\mathbb S/\eta^4=\mathbb S\oplus\Sigma^5\mathbb S$, which contains $\mathbb S$ as a retract. Thus, we get all perfect $\mathbb S$-modules! Note that this example suggests that we will have to use some machinery: it is not obvious that all elements in $\pi_*\mathbb S$ are nilpotent!
\end{example}
\begin{proof}[Proof of \Cref{thm:thick-finite-subcat}]
	It is not too hard to check that $\mc C_{\ge n}$ is actually thick. This is a matter of checking that the type can only increase when taking various finite (co)limits and retracts and so on.

	Thus, the hard part is to show the following: if $\mc C$ is a thick subcategory containing a finite spectrum of type $n$, then $\mc C$ contains every finite spectrum of type at least $n$. Before proceeding with the argument, we establish some notation.
	\begin{definition}[dual]
		For a finite spectrum $X$, we define the \textit{dual} to be
		\[DX\coloneqq\op{Hom}_{\mathrm{Spectra}}(X,\mathbb S_{(p)}).\]
	\end{definition}
	\begin{remark}
		Note that $DX$ is a finite spectrum: certainly $D\mathbb S_{(p)}=\mathbb S_{(p)}$ is a finite spectrum, and general $X$ can be built out of retracts of finite (co)limits, which commute with $\op{Hom}$.
	\end{remark}
	\begin{remark}
		Using similar arguments to the previous remark, one can check that the canonical map
		\[X_1\otimes DX_2\to\op{Hom}_{\mathrm{Spectra}}(X_2,X_1)\]
		is an isomorphism.
	\end{remark}
	\begin{notation}
		For a finite spectrum $X$, we also define $\op{End}X\coloneqq X\otimes DX$.
	\end{notation}
	\begin{remark}
		We note that $\op{End}X$ is an $\mathbb E_1$-ring.
	\end{remark}
	We now proceed with the argument.
	\begin{enumerate}
		\item We now proceed with the proof. Fix our $X\in\mc C$ of type $n$. Then $\op{End}X\coloneqq X\otimes DX$ is an $\mathbb E_1$-ring, so it admits a unit $e\colon\mathbb S_{(p)}\to\op{End}X$, whose fiber we denote by $f\colon F\to\mathbb S_{(p)}$. Thus, we have built a fiber sequence
		\[F\stackrel f\to\mathbb S_{(p)}\stackrel e\to\op{End}X.\]
		Because $K(m)_*X\ne0$ for $m\ge n$, we see that $K(m)_*e$ is injective: this amounts to asserting that the target ring $K(m)_*\op{End}X$ is nonzero! The long exact sequence thus implies that $K(m)_*f=0$ for $m\ge n$.
		\item Now, choose some $Y$ of type at least $n$, and we would like to check that $Y$ is in $\mc C$. We still have a unit map $\mathbb S_{(p)}\to Y\otimes DY$, so the composite
		\[F\to\mathbb S_{(p)}\to\op{End}Y\]
		vanishes on $K(m)_*$ for $m\ge n$ still by the previous paragraph. Note that this implies that the composite vanishes on $\FF_{p*}$ as well: $\FF_p$ is some infinite colimit of the $K(m)$s, so because $Y$ is finite, the $\FF_p$-homology will agree with $K(m)$-homology for $m$ large enough.\footnote{This is a little tricky: imagine taking $-\otimes K(m)$, and because $K(m)$ is a field, everything splits, and one can try hard to compute with a spectral sequence.} We also remark that the composite vanishes on $K(m)_*$ for $m<n$ because $K(m)_*Y=0$ for $m<n$ already.
		\item We now use our machinery. By \Cref{cor:map-to-bp-null} and \Cref{prop:bp-bousfield}, we conclude that there is a $k$ large enough for which the composite
		\[F^{\otimes k}\to(\op{End}Y)^{\otimes k}\to\op{End}Y\]
		vanishes. By moving the duals around, we can show that
		\[F^{\otimes k}\otimes Y\stackrel f\to F^{\otimes(k-1)}\otimes Y\stackrel f\to\cdots\stackrel f\to Y\]
		is trivial for some $k$.
		\item We complete the proof. By the previous step, we conclude that $Y$ is a retract of $Y/(F^{\otimes k}\otimes Y)$ by the long exact sequence in homotopy. Splitting up this large tensor product, we see that $Y$ is in the thick subcategory generated by $(F^{\otimes a}\otimes Y)/(F^{\otimes(a+1)}\otimes Y)$s, but this latter object is $F^{\otimes a}\otimes Y\otimes DX$. This last object is a tensor product of $X$ with a finite spectrum, which is in $\mc C$: indeed, tensor products commute with finite limits, colimits, and retracts, so this tensor product can be turned into a retract of some finite (co)limit of $X$.
		\qedhere
	\end{enumerate}
\end{proof}

\subsection{Self Maps of Finite Spectra}
The following example motivates our next theorem.
\begin{example}
	Note that $\mathbb S$ has type $0$, and it has a self map $p\colon\mathbb S\to\mathbb S$ with cofiber of type $1$. Similarly, $\mathbb S/p$ is type $1$, and it has a self map with cofiber of type $2$: for $p>2$, we can take $v_1\colon\mathbb S^{2p-2}/p\to\mathbb S/p$; for $p=2$, we can take $v_1^4\colon\mathbb S^8/2\to\mathbb S/2$.
\end{example}
Let's codify a good property of these self-maps.
\begin{definition}
	Fix a finite spectrum $X$. Then a \textit{$v_n$-self map} is an endomorphism $f\colon\Sigma^KX\to X$ such that $K(n)_*f$ is an isomorphism, but $K(m)_*f$ is nilpotent for $m\ne n$.
\end{definition}
These maps all exist.
\begin{theorem} \label{thm:vn-maps-exist}
	Let $X$ be a finite spectrum of type $n$. Then $X$ admits a $v_n$-self map $\Sigma^kX\to X$.
\end{theorem}
\begin{remark}
	This is intended to be a little similar to our story with slopes: given a particular slope, we have a self-map, and taking the quotient by a self-map should give a smaller slope.
\end{remark}
\begin{remark}
	Note that the cofiber of a $v_n$-self map automatically has type $n+1$, which can be seen by taking $K(n)_*$ of the corresponding fiber sequence.
\end{remark}
Before proceeding with the proof, we need a bit more language.
\begin{definition}
	Fix a finite $p$-local ring $\mathbb E_1$-ring $R$. Then a \textit{$v_n$-element} in $\pi_*R$ is a homotopy class which maps to a unit in $K(n)_*R$ and a nilpotent element in $K(m)_*R$ for $m\ne n$.
\end{definition}
\begin{remark}
	By unwinding the definitions, we see that a $v_n$-self map of $X$ is a $v_n$-element in $\pi_*(X\otimes DX)$.
\end{remark}
The following lemma explains the terminology.
\begin{lemma} \label{lem:high-vn-power-behaves}
	Let $x$ be a $v_n$-element of a finite $p$-local ring $R$. Then there is $N$ such that $K(m)_*x^N=0$ for $m\ne n$ and such that $K(n)_*x$ is a power of $v_n$.
\end{lemma}
\begin{proof}
	Everything is perfect, so the homology is given by just $\FF_{p*}$ for $m$ large enough, so we just choose $N$ large enough so that $K(m)_*x^N=0$ for $m$ small and $\FF_{p*}x^N=0$. As for the second requirement, note that $K(n)_*R/(v_n-1)$ is a finite set (we are kind of looking at the group of units of a finite ring), so this is just a congruence condition on $N$.
\end{proof}
These $v_n$-elements turn out to have a lot of structure. To begin, we will also want the following piece of algebra.
\begin{lemma} \label{lem:get-powers-equal}
	Let $A$ be a discrete associative ring. Suppose that $x,y\in A$ are chosen such that $xy=yx$ and such that $x-y$ is $p$-power torsion and nilpotent. Then $x^{p^N}=y^{p^N}$ for $N$ large enough.
\end{lemma}
\begin{proof}
	Write $x^{p^N}$ as
	\[(y+(x-y))^{p^N},\]
	and we imagine expanding this out using the binomial theorem. There is always a ``main term'' $y^{p^N}$, and the remaining terms look like $\binom{p^N}dy^{p^N-d}(x-y)^d$. For larger $d$, this automatically vanishes by nilpotence, and for large enough $N$, we can make the smaller terms have $\binom{p^N}d$ divisible by a high power of $p$.
\end{proof}
Here is our application of the algebra.
\begin{lemma} \label{lem:high-power-commutes}
	Let $R$ be a finite $p$-local ring, and suppose that $x$ is a $v_n$-element in $\pi_*R$ (for some positive $n$). Then some power of $x$ is central in $\pi_*R$.
\end{lemma}
\begin{proof}
	We consider the ring $R_2\coloneqq X\otimes DX$. Now, $R_2$ has two elements $a$ and $b$ which are given by left- and right-multiplication by $x$. By \Cref{lem:get-powers-equal}, it is enough to check that $a$ and $b$ commute, have difference which is torsion, and $a-b$ is nilpotent. The commutativity has little content (because left- and right-multiplication should always commute), and being torsion follows because $R$ is finite (namely, everything has trivial homology over $\QQ$). Lastly, $a-b$ is nilpotent by \Cref{lem:high-vn-power-behaves}: it is enough to check the vanishing for every $K(n)$, and the lemma tells us that a high power of $a$ and $b$ acts the same on both sides, so the nilpotence basically follows.
\end{proof}
\begin{proposition} \label{prop:vn-powers}
	Let $R$ be a finite $p$-local ring, and suppose that $x$ and $y$ are $v_n$-elements. Then $x^a=y^b$ for some $a$ and $b$.
\end{proposition}
\begin{proof}
	By \Cref{lem:high-power-commutes}, we may replace $x$ and $y$ with powers which commute. Then the same argument as in \Cref{lem:high-power-commutes} shows that a high enough power of $x$ and $y$ must be equal.
\end{proof}
\begin{corollary} \label{cor:vn-powers-functorial}
	Let $f\colon X\to Y$ be a morphism of finite spectra. If $X$ and $Y$ admit $v_n$-self maps $g$ and $h$, then there is a commutative diagram
	% https://q.uiver.app/#q=WzAsNCxbMCwwLCJYIl0sWzEsMCwiWSJdLFswLDEsIlxcU2lnbWFea1giXSxbMSwxLCJcXFNpZ21hXmtZIl0sWzAsMSwiZiJdLFswLDIsImdeYSIsMl0sWzEsMywiaF5iIl0sWzIsMywiZiJdXQ==&macro_url=https%3A%2F%2Fraw.githubusercontent.com%2FdFoiler%2Fnotes%2Fmaster%2Fnir.tex
	\[\begin{tikzcd}[cramped]
		X & Y \\
		{\Sigma^kX} & {\Sigma^kY}
		\arrow["f", from=1-1, to=1-2]
		\arrow["{g^a}"', from=1-1, to=2-1]
		\arrow["{h^b}", from=1-2, to=2-2]
		\arrow["f", from=2-1, to=2-2]
	\end{tikzcd}\]
	for some $a$ and $b$.
\end{corollary}
\begin{proof}
	Both $h$ and $Dg$ are $v_n$-self maps of $Y\otimes DX$, so some power of them agree by \Cref{prop:vn-powers}. The given equality follows by evaluating $1$ on the composite
	\[\mathbb S\stackrel f\to Y\otimes DX\to Y\otimes DX,\]
	where the second map is $h^a=Dg^b$.
\end{proof}
We are now ready to prove \Cref{thm:vn-maps-exist}. We will instead prove the following stronger result.
\begin{theorem}
	Let $X$ be a finite spectrum of type $n$. Then $X$ admits a $v_n$-self map, and it is unique up to taking powers.
\end{theorem}
\begin{proof}
	The uniqueness was already shown in \Cref{prop:vn-powers}. By looking at some explicit examples of finite spectra of type $n$ (e.g., from lens spaces), it only remains to show that the subcategory of spectra of type at least $n$ admitting a $v_n$-self map is a thick subcategory. It's not hard to show that admitting such a map is preserved by produces or retracts, so it only remains to show the claim for cofibers. But this follows by taking a cofiber of $f\colon X\to Y$ in the diagram in \Cref{cor:vn-powers-functorial}.
\end{proof}

\subsection{Bousfield Localization}
Our next goal is to complete a spectrum. To this end, we will discuss the more general notion of Bousfield localization. We will need a size condition.
\begin{definition}[presentable]
	Fix a ring spectrum $R$. Then a full subcategory $\mc C\subseteq\mathrm{Mod}(R)$ is \textit{presentable} if and only if it satisfies the following.
	\begin{itemize}
		\item Closure: $\mc C$ is closed under de-suspensions and all colimits.
		\item Size: there is a set of objects generating $\mc C$ under colimits.
	\end{itemize}
\end{definition}
Here is the main theorem; it is an application of some Adjoint functor theorem.
\begin{theorem} \label{thm:bousfield-localization}
	Let $R$ be a ring spectrum, and let $\mc C\subseteq\mathrm{Mod}(R)$ be a presentable full subcategory. Then the inclusion $\mc C\to\mathrm{Mod}(R)$ preserves colimits and admits a right adjoint $G$ given by
	\[G(M)\coloneqq\colim_{\substack{N\in\mc C\\f\colon N\to M}}N.\]
\end{theorem}
\begin{remark}
	We have tended to ignore set-theoretic conditions, but this one is important because we are trying to apply an Adjoint functor theorem. In particular, the colimit does not even really make sense if $\mc C$ is not even generated by a set of objects.
\end{remark}
\Cref{thm:bousfield-localization} gives rise to our localization.
\begin{definition}[Bousfield localization]
	Let $\mc C$ be a presentable full subcategory of $\mathrm{Mod}(R)$, where $R$ is a ring spectrum, and let $G_{\mc C}$ be the right adjoint of the inclusion. For an $R$-module $M$, the cofiber of the unit map $G_{\mc C}M\to M$ is the \textit{Bousfield localization}, labeled $L_{\mc C}M$.
\end{definition}
\begin{example}
	Let $\mc C$ denote the collection of $\mathbb S_{(p)}$-modules generated under all colimits from $\mathbb S/p$ and its de-suspensions. Then what is left turns out to be $L_{\mc C}X=X[1/p]$. Thus, Bousfield localization recovers usual localization.
\end{example}
\begin{remark}
	In fact, we can use \Cref{thm:bousfield-localization} to show that $R[1/r]$ is an $\mathbb E_\infty$-ring for any $r\in\pi_*R$, essentially by identifying its category of modules.
\end{remark}
Remarkably, we can also recover completion.
\begin{example}
	Let $\mc C$ denote the collection of $\mathbb S_{(p)}$-modules for which the canonical map $M\to M[1/p]$. (It is not obvious that this subcategory is generated by a set of objects!) Then $L_{\mc C}X$ is the completion $X^\land_p$.
\end{example}

\end{document}